\documentclass[10pt,a4paper]{amsart}
\usepackage[margin=19mm]{geometry}
\usepackage{amsmath,amssymb,mathtools}
\usepackage[scr=rsfs,cal=euler]{mathalfa}
\usepackage{xfrac}
\usepackage[unicode,hidelinks,bookmarks=false]{hyperref}
\newcommand{\ie}{i.e.\ }
\newcommand{\cf}{cf.\ }
\newcommand{\resp}{respectively\ }
\newcommand{\Subsection}{\subsection}
\providecommand{\nobreakdash}{\nobreak\hbox{-}}
\setlength{\emergencystretch}{2em}
\multlinegap0pt
\usepackage{bm}
\usepackage{bbm}
\usepackage{dsfont}
\usepackage{xspace}
\usepackage{cancel}
\usepackage{mathtools}
\usepackage{amsthm}
\makeatletter
\@ifundefined{thm}{\newtheorem{thm}{Theorem}}{}
\@ifundefined{defn}{\newtheorem{defn}[thm]{Definition}}{}
\@ifundefined{lem}{\newtheorem{lem}[thm]{Lemma}}{}
\@ifundefined{prop}{\newtheorem{prop}[thm]{Proposition}}{}
\makeatother
\newcommand{\CO}{\mathcal{O}}         
\newcommand{\Orb}{\operatorname{Orb}} 
\newcommand{\mb}[1]{\mathbf{#1}}      
\title[A Proof of the Biquadratic Linear AFL for GL(4)]{A Proof of the Biquadratic Linear AFL for GL(4)}
\author{Qirui Li}
\date{Source: arXiv:2505.22625v1 (CC BY 4.0)}
\begin{document}
\maketitle

\noindent\emph{Source and licence.} A Proof of the Biquadratic Linear AFL for GL(4). Version arXiv:2505.22625v1 (CC BY 4.0), distributed under the Creative Commons Attribution 4.0 International licence, as recorded in the arXiv licence metadata for this exact version and shown on the abstract page https://arxiv.org/abs/2505.22625v1. Full rights notice: licenses/arxiv-2505.22625-CC-BY-4.0.txt.\par\smallskip

\noindent\textit{Editorial scope.} Counted unit: the Thm whose source label is \texttt{orbitalh=2} in the article (source0.tex lines 1570--1588, source label \texttt{orbitalh=2}), delivered together with the article's own complete original proof of it (source0.tex lines 1590--1675, one \texttt{proof} environment of 86 lines). No mathematics is written, repaired or re-derived: statement and proof are the article's own text line for line. Required context retained uncounted: index (lines 1324--1330); orbits (lines 1332--1342); pilemma (lines 1520--1525). Every definition, display and label of those lines is retained unchanged. Omitted from this package and not redistributed: 1--1323, 1331--1331, 1343--1519, 1526--1569, 1589--1589, 1676--2434. Nothing inside a retained excerpt is omitted or altered: every retained body is delivered exactly as the article's lines are written, so no omission receipt of any kind accompanies this package. Citations: the retained excerpts carry no citation command at all, so no bibliography entry of the article is transcribed and no reference is redistributed. Reference adaptation: none was needed and none was made; every reference inside the delivered unit resolves inside it, so no unresolved reference or citation is produced. Printed slips kept literal: no printed word, symbol, hypothesis or number of the counted statement or of its proof was altered. Layout and class: the article's own class is \texttt{article} with packages including amsmath, amssymb, amsthm, mathrsfs, xy, inputenc, fontenc, geometry, xr-hyper, hyperref, color; the wrapper is the lane's pinned \texttt{amsart} editorial wrapper at 10pt on A4, which restates the theorem-like environments the retained text needs and adds the article's own macro definitions that the retained text needs (\texttt{\textbackslash CO}, \texttt{\textbackslash Orb}, \texttt{\textbackslash mb}), with extra wrapper packages (bm, bbm, dsfont, xspace, cancel, mathtools). The delivered numbering is the unit's own. Assets: no retained span contains a figure environment, a graphics-inclusion command, a PDF-page inclusion, a table or a TikZ picture; no image file is copied into this package, the package declares no asset dependency and no image file is redistributed with it. Licence: version arXiv:2505.22625v1 of this article is distributed under the Creative Commons Attribution 4.0 International licence, as recorded in the arXiv licence metadata for this exact version and shown on the abstract page https://arxiv.org/abs/2505.22625v1. A full rights notice is delivered at licenses/arxiv-2505.22625-CC-BY-4.0.txt. Characters: every retained excerpt is pure ASCII, so the delivered engine, XeLaTeX with its default Latin Modern text font, reports no missing character.
\begin{defn}\label{index}
We write
$$
\Lambda_1 \overset{k}{\subset} \Lambda_2
$$
to indicate that $\Lambda_1 \subset \Lambda_2$ and $\#(\Lambda_2 / \Lambda_1) = q^k$.
\end{defn}

\begin{prop}\label{orbits}
For any $n \geq 1$, there are $q$ sublattices
$$
\Lambda \overset{1}{\subset} R_n
$$
such that $\Lambda \cong R_{n+1}$, and there is a unique lattice
$$
\Lambda \overset{1}{\subset} R_n
$$
such that $\Lambda \cong R_{n-1}$.
\end{prop}

\begin{lem}\label{pilemma}
If $|\mb w|_L < |\varpi_3|_L$, then for any $n \leq r$, we have
$$
\mb z^2 \cdot R_n = \pi R_n.
$$
\end{lem}

\begin{thm}\label{orbitalh=2}
If $L/F$ is an unramified extension, then
$$
\Orb(\mathbf{1}, \beta, s) = (1 + q^{2s}) + (1 - q^s)^2 \cdot (1 + q^{-1}) \cdot (q + q^2 + \cdots + q^r).
$$
In particular,
\begin{equation}\label{unramifiedcen}
\Orb(\mathbf{1}, \beta, 0) = 2.
\end{equation}

If $L/F$ is a ramified extension, then
$$
\Orb(\mathbf{1}, \beta, s) = (1 - q^s + q^{2s}) + (1 - q^s)^2 \cdot (q + q^2 + \cdots + q^r).
$$
In particular,
\begin{equation}\label{ramifiedcen}
\Orb(\mathbf{1}, \beta, 0) = 2.
\end{equation}
\end{thm}

\begin{proof}
Recall Definition~\ref{index}. We compute the orbital integral by collecting coefficients of $(-q^s)^k$. Since $|\mb z^2|_L = q^2$, we may write
$$
\Orb(\mathbf{1}, \beta, s) = a_0 + a_1(-q^s) + a_2(-q^s)^2,
$$
where
\begin{align*}
a_0 &= \sum_{n=0}^r \left[\CO_L^\times : R_n^\times\right] \cdot \#\left\{ \Lambda : R_n = \Lambda \supset \mb z^2 \cdot R_n \right\}, \\
a_1 &= \sum_{n=0}^r \left[\CO_L^\times : R_n^\times\right] \cdot \#\left\{ \Lambda : R_n \overset{1}{\supset} \Lambda \overset{1}{\supset} \mb z^2 \cdot R_n,\ \Lambda \cong R_k,\ k \leq r \right\}, \\
a_2 &= \sum_{n=0}^r \left[\CO_L^\times : R_n^\times\right] \cdot \#\left\{ \Lambda : R_n \supset \Lambda = \mb z^2 \cdot R_n \right\}.
\end{align*}

Clearly,
\begin{equation}\label{a02}
a_0 = a_2 = \sum_{n=0}^r \left[\CO_L^\times : R_n^\times\right].
\end{equation}

To compute $a_1$, we apply Lemma~\ref{pilemma}, and decompose
$$
a_1 = I_0 + \sum_{n=1}^r J_n + \sum_{n=1}^r K_n,
$$
where
\begin{align*}
I_0 &= \left[\CO_L^\times : \CO_L^\times\right] \cdot \#\left\{ \Lambda \cong \CO_L : \CO_L \overset{1}{\supset} \Lambda \overset{1}{\supset} \pi \CO_L \right\}, \\
J_n &= \left[\CO_L^\times : R_{n-1}^\times\right] \cdot \#\left\{ \Lambda \cong R_n : R_{n-1} \overset{1}{\supset} \Lambda \overset{1}{\supset} \pi R_{n-1} \right\}, \\
K_n &= \left[\CO_L^\times : R_n^\times\right] \cdot \#\left\{ \Lambda \cong R_{n-1} : R_n \overset{1}{\supset} \Lambda \overset{1}{\supset} \pi R_n \right\}.
\end{align*}

By Proposition~\ref{orbits}, we have:
$$
J_n = \left[\CO_L^\times : R_n^\times\right], \qquad
K_n = \left[\CO_L^\times : R_n^\times\right].
$$

For $I_0$, we distinguish cases:
$$
I_0 =
\begin{cases}
0 & \text{if $L/F$ is unramified}, \\
1 & \text{if $L/F$ is ramified}.
\end{cases}
$$

Thus:
\begin{itemize}
\item If $L/F$ is unramified:
$$
a_1 = 0 + 2 \sum_{n=1}^r \left[\CO_L^\times : R_n^\times\right] = 2(a_0 - 1);
$$
\item If $L/F$ is ramified:
$$
a_1 = 1 + 2 \sum_{n=1}^r \left[\CO_L^\times : R_n^\times\right] = 1 + 2(a_0 - 1).
$$
\end{itemize}

Hence we may write:
$$
a_1 = c_L + 2(a_0 - 1), \qquad \text{where} \quad
c_L =
\begin{cases}
0 & \text{if $L/F$ unramified}, \\
1 & \text{if $L/F$ ramified}.
\end{cases}
$$

Substituting into the orbital integral expression:
\begin{align*}
\Orb(\mathbf{1}, \beta, s)
&= a_0 + a_1(-q^s) + a_2(-q^s)^2 \\
&= (a_0 - 1)(1 - q^s)^2 + \left(1 + c_L \cdot (-q)^s + (-q^s)^2\right).
\end{align*}

To finish, we evaluate $a_0 - 1$ in each case:
\begin{itemize}
\item If $L/F$ is unramified:
$$
a_0 - 1 = \sum_{n=1}^r \left[\CO_L^\times : R_n^\times\right] = \sum_{n=1}^r (1 + q^{-1}) q^n;
$$
\item If $L/F$ is ramified:
$$
a_0 - 1 = \sum_{n=1}^r \left[\CO_L^\times : R_n^\times\right] = \sum_{n=1}^r q^n.
$$
\end{itemize}

This concludes the proof.
\end{proof}

\end{document}
